\documentclass[11pt]{article}
\usepackage[margin=1in]{geometry}
\usepackage{booktabs}
\usepackage{hyperref}
\usepackage{xcolor}
\usepackage{listings}
\usepackage{tabularx}

\hypersetup{colorlinks=true, linkcolor=blue!60!black, urlcolor=blue!60!black}
\lstset{basicstyle=\ttfamily\small, breaklines=true, frame=single,
        backgroundcolor=\color{gray!8}}

\title{Interference-Aware Routing:\\Evaluation Report}
\author{Autonomous Networks Research Group (ANRG)\\University of Southern California}
\date{}

\begin{document}
\maketitle

%======================================================================
\section{Experiment Setup}

\textbf{Simulator:} ncsim, HEFT scheduler, \texttt{csma\_bianchi} 802.11ax interference model.
\textbf{RF config:} 5\,GHz, 20\,MHz channel, $P_{\mathrm{tx}}=20$\,dBm, path-loss exponent $n=3.0$.
\textbf{Grid spacing:} 40\,m (adjacent links $\approx8.6$\,MB/s PHY; diagonals $\approx56.6$\,m, lower PHY rate).
\textbf{Seeds:} Each configuration averaged over 30 seeds (smoothing HEFT non-determinism from Python hash randomization).
\textbf{Workload:} Heterogeneous compute costs (150--1000) and data sizes (2--30\,MB).

\subsection{Networks}

\begin{table}[h]
\centering
\begin{tabular}{lrrp{6.5cm}}
\toprule
\textbf{Network} & \textbf{Nodes} & \textbf{Links} & \textbf{Topology} \\
\midrule
$4\times4$ grid & 16 & 68  & Grid + checkerboard diagonals, heterogeneous compute \\
$7\times7$ grid & 49 & 240 & Grid + checkerboard diagonals, heterogeneous compute \\
\bottomrule
\end{tabular}
\end{table}

\subsection{DAGs}

\begin{table}[h]
\centering
\begin{tabular}{lrrp{6cm}}
\toprule
\textbf{DAG} & \textbf{Tasks} & \textbf{Edges} & \textbf{Structure} \\
\midrule
Small           &  8 &  12 & Fork-join: 1 source $\to$ 6 parallel $\to$ 1 sink \\
Large ($4\times4$) & 30 &  48 & 5-stage pipeline with cross-connections \\
Large ($7\times7$) & 60 & 102 & 6-stage pipeline with cross-connections \\
\bottomrule
\end{tabular}
\end{table}

\subsection{Schemes Evaluated}

\begin{table}[h]
\centering
\small
\begin{tabular}{llp{7.5cm}}
\toprule
\textbf{Label} & \textbf{Scheme} & \textbf{Description} \\
\midrule
W     & Widest-Path        & Max bottleneck bandwidth; ignores interference \\
S     & Shortest-Path      & Min total latency ($\equiv$ min hops at equal latency); ignores interference \\
GS    & Greedy-Start       & Static greedy at DAG inject; flows sorted by est.\ start time \\
GC    & Greedy-Criticality & Static greedy; flows sorted by HEFT upward rank (desc.) \\
GB    & Greedy-Bytes       & Static greedy; flows sorted by data size (desc.) \\
GO    & Greedy-Overlap     & Static greedy; flows sorted by overlap degree (desc.) \\
GSD   & Dynamic            & Routes computed at \texttt{transfer\_start} using actual active-link state \\
GSD-D & Dynamic + Deferral & Like GSD, but defers transfers when congestion exceeds threshold \\
\bottomrule
\end{tabular}
\end{table}

%======================================================================
\section{Results}

\subsection{Mean Makespan (seconds, averaged over 30 seeds)}

\begin{table}[h]
\centering
\small
\begin{tabular}{lrrrrrrrrl}
\toprule
\textbf{Experiment}
  & \textbf{W} & \textbf{S} & \textbf{GS} & \textbf{GC}
  & \textbf{GB} & \textbf{GO} & \textbf{GSD} & \textbf{GSD-D} & \textbf{Best} \\
\midrule
$4\times4$, small &  191.1 &  \textbf{20.4} &   50.9 &   50.9 &   50.9 &   50.9 &   76.0 &   48.3 & S     \\
$4\times4$, large &  363.5 &  234.1 &  183.8 &  203.4 &  178.4 &  193.6 &  298.7 & \textbf{158.4} & GSD-D \\
$7\times7$, small &  250.9 &  \textbf{35.8} &  187.8 &  184.6 &  202.4 &  192.8 &  230.2 &  199.2 & S     \\
$7\times7$, large & 4852.9 & \textbf{2465.1} & 3644.7 & 3850.1 & 3518.1 & 3873.5 & 3603.5 & 3629.5 & S     \\
\bottomrule
\end{tabular}
\end{table}

\subsection{Improvement over Widest-Path Baseline}

\begin{table}[h]
\centering
\begin{tabular}{lrrrrrrr}
\toprule
\textbf{Experiment}
  & \textbf{S} & \textbf{GS} & \textbf{GC} & \textbf{GB} & \textbf{GO} & \textbf{GSD} & \textbf{GSD-D} \\
\midrule
$4\times4$, small & $\mathbf{-89\%}$ & $-73\%$ & $-73\%$ & $-73\%$ & $-73\%$ & $-60\%$ & $-75\%$ \\
$4\times4$, large & $-36\%$ & $-49\%$ & $-44\%$ & $-51\%$ & $-47\%$ & $-18\%$ & $\mathbf{-56\%}$ \\
$7\times7$, small & $\mathbf{-86\%}$ & $-25\%$ & $-26\%$ & $-19\%$ & $-23\%$ & $-8\%$  & $-21\%$ \\
$7\times7$, large & $\mathbf{-49\%}$ & $-25\%$ & $-21\%$ & $-28\%$ & $-20\%$ & $-26\%$ & $-25\%$ \\
\bottomrule
\end{tabular}
\end{table}

\subsection{Winners}

\begin{table}[h]
\centering
\begin{tabular}{llll}
\toprule
\textbf{Experiment} & \textbf{Winner} & \textbf{Runner-up} & \textbf{Gap} \\
\midrule
$4\times4$, small & \textbf{S} (20.4\,s)     & GSD-D (48.3\,s)  & 57.7\% \\
$4\times4$, large & \textbf{GSD-D} (158.4\,s) & GB (178.4\,s)    & 11.2\% \\
$7\times7$, small & \textbf{S} (35.8\,s)      & GC (184.6\,s)    & 80.6\% \\
$7\times7$, large & \textbf{S} (2465.1\,s)    & GB (3518.1\,s)   & 29.9\% \\
\bottomrule
\end{tabular}
\end{table}

%======================================================================
\section{Analysis}

\subsection{Shortest-Path Dominates with Heterogeneous Workloads}

S wins 3 of 4 experiments.  With heterogeneous data sizes (2--30\,MB),
min-hop paths are especially valuable for large transfers.  A 30\,MB
transfer on a 1-hop path avoids the multi-hop relay chains that trigger
severe \texttt{csma\_bianchi} interference.  S's advantage is largest on
small DAGs ($-86$--$89\%$ vs W) where contention is low enough that simply
avoiding unnecessary hops is optimal.

\subsection{GSD-D Wins on $4\times4$ with Large DAGs}

On the $4\times4$ grid with 30 tasks, GSD-D (158.4\,s) is the clear winner,
beating GB (178.4\,s) by 11\% and S (234.1\,s) by 32\%.  This is the
\textbf{only scenario where deferral helps}: the small, dense grid forces
many transfers onto the same few links, and deferral allows transfers to
wait for less-congested conditions.

Notably, GSD-D dramatically outperforms its non-deferral counterpart GSD
(298.7\,s vs 158.4\,s, $-47\%$).  Without deferral, GSD's per-flow greedy
optimization makes locally optimal but globally suboptimal choices on the
dense grid.  Deferral breaks this cascade by letting the engine wait for
airtime to free up.

\subsection{Deferral Does Not Help on $7\times7$ Grids}

On the $7\times7$ grid, GSD-D (199.2\,s small, 3629.5\,s large) provides no
improvement over the static greedy variants.  The larger grid offers enough
alternative paths that transfers can route around congestion without
deferring.  Deferral adds latency without commensurate benefit.

\subsection{Greedy Orderings Diverge with Heterogeneous Data}

With heterogeneous data sizes, the greedy orderings now show meaningful
differentiation.  On $4\times4$ large, GB (178.4\,s) beats GS (183.8\,s) by
3\%---GB's largest-first ordering gives the 30\,MB transfers the best
routes.  On $7\times7$ large, GB (3518.1\,s) outperforms all other greedy
variants by 2--10\%.

The four static greedy variants are no longer tied on $4\times4$ small
(50.93\,s each).  With uniform data sizes they were identical; with
heterogeneous sizes, the ordering happens to produce the same routes for
this simple fork-join topology.

\subsection{Widest-Path Remains Consistently the Worst}

W is $5$--$9\times$ worse than the best scheme.  With heterogeneous data
sizes, the penalty for long relay chains is even more severe: a 30\,MB
transfer on a 3-hop Widest-Path route experiences compounding interference
on each link.

%======================================================================
\section{Regime Map}

\begin{table}[h]
\centering
\begin{tabular}{lll}
\toprule
\textbf{Scenario} & \textbf{Best approach} & \textbf{Why} \\
\midrule
Small grid, large DAG    & \textbf{GSD-D} & Dense contention; deferral reduces interference \\
Small grid, small DAG    & \textbf{S}     & Few transfers; min-hop avoids interference \\
Large grid, any DAG      & \textbf{S}     & Enough paths to spread traffic; min-hop is sufficient \\
Heterogeneous data sizes & \textbf{GB}    & Largest-first ordering gives big transfers best routes \\
\bottomrule
\end{tabular}
\end{table}

%======================================================================
\section{Deferral Analysis}

\subsection{When Deferral Helps}

GSD-D's deferral mechanism (threshold=0.3, max\_deferrals=3) is effective
only on the $4\times4$ grid with a large DAG.  This is the regime with the
highest link contention: 48 data edges compete for 68 links on a 16-node
grid.  The deferral mechanism allows transfers to wait 1--2 transfer
completions before starting, finding less-congested paths.

\subsection{When Deferral Hurts}

On the $7\times7$ grid, deferral adds no value and slightly hurts
performance compared to static greedy schemes.  The larger grid provides
enough routing diversity that transfers can always find an acceptable path
immediately.  Deferring only delays the transfer without improving the path
quality on retry.

\subsection{Comparison: GSD vs GSD-D}

\begin{table}[h]
\centering
\begin{tabular}{lrrr}
\toprule
\textbf{Experiment} & \textbf{GSD} & \textbf{GSD-D} & \textbf{Improvement} \\
\midrule
$4\times4$, small &   76.0 &   48.3 & $-36\%$ \\
$4\times4$, large &  298.7 &  158.4 & $-47\%$ \\
$7\times7$, small &  230.2 &  199.2 & $-13\%$ \\
$7\times7$, large & 3603.5 & 3629.5 & $+1\%$  \\
\bottomrule
\end{tabular}
\end{table}

Deferral consistently improves on GSD ($-13$ to $-47\%$) except on
$7\times7$ large where it is essentially neutral ($+1\%$).  The largest
gain is on $4\times4$ large ($-47\%$), confirming that deferral is most
valuable in high-contention regimes.

%======================================================================
\section{Reproducing These Results}

\begin{lstlisting}
cd ncsim/
python run_routing_eval.py
# Output: /tmp/ncsim_routing_eval/eval_results.json
# Runtime: ~50 minutes (960 runs)
\end{lstlisting}

\end{document}
